\documentclass[12pt, letterpaper]{article}

%% Build from the paper/ directory:
%%   cd paper && pdflatex main.tex
%% or use the root build system:
%%   make paper          (from repo root)
%%   ./build.sh paper    (POSIX)
%%   build.cmd paper     (Windows cmd)

% packages.tex
\usepackage{amsmath, amssymb, amsthm}
\usepackage{graphicx}
\usepackage{xcolor}
\usepackage[
    hidelinks,
    pdftitle={Foundations of the A=1 Discrete Causal Lattice: Bipartite Octahedral Lattices in Arbitrary Dimension},
    pdfauthor={Jack Menendez},
    pdfsubject={Foundations mathematics for the A=1 Discrete Causal Lattice series},
    pdfkeywords={discrete causal lattice, bipartite octahedral lattice, discrete differential geometry, discrete-to-continuum limit, Hilbert's sixth problem, foundations of physics}
]{hyperref}
\usepackage{booktabs}
\usepackage{array}
\usepackage{geometry}
\usepackage{adjustbox}
\usepackage{longtable}
\usepackage{setspace}
\usepackage{caption}
\usepackage{subcaption}
\usepackage{float}
\onehalfspacing

\input{macros/commands}

%% Page geometry (also set here as a fallback in case the package path fails)
\geometry{top=1in, bottom=1in, left=1in, right=1in}

%% ---------------------------------------------------------------
%% TITLE PAGE
%% ---------------------------------------------------------------

\title{Foundations of the A=1 Discrete Causal Lattice:\\
       the Diamond Progression in Arbitrary Dimension%
       \thanks{Version 0.2-DRAFT.
       \href{https://github.com/JackDMenendez/dcl-mathematics}%
            {GitHub: JackDMenendez/dcl-mathematics}.
       %% Fill in at release time (conditional on publication decision):
       %% \href{https://doi.org/10.5281/zenodo.XXXXXXXX}%
       %%      {DOI: 10.5281/zenodo.XXXXXXXX}.
       Pre-release draft; DOI to be assigned if and when an external
       publication decision is taken.}}
\author{Jack Menendez\\
        \small Independent Scholar, Eugene, OR\\
        \small ORCID: \href{https://orcid.org/0009-0003-1166-307X}{0009-0003-1166-307X}\\
        \small \href{mailto:jackdmenendez@gmail.com}{jackdmenendez@gmail.com}}
\date{\today}

\begin{document}

\maketitle

%% ---------------------------------------------------------------
%% FRONT-MATTER NOTICES
%% ---------------------------------------------------------------
\noindent\textbf{License.}  Paper text and figures: Creative Commons
Attribution 4.0 (CC BY 4.0).  Source code in the companion
repository: MIT License.

\medskip
\noindent\textbf{Data availability.}  The LaTeX source for this paper
is publicly available at
\href{https://github.com/JackDMenendez/dcl-mathematics}%
     {\nolinkurl{github.com/JackDMenendez/dcl-mathematics}}.
This is a paper-only repository in the A=1 Discrete Causal Lattice
(DCL) series; no numerical experiments accompany the first paper.
%% At release time (if external publication proceeds):
%% The exact commit corresponding to this release is the vX.Y tag,
%% archived at
%% \href{https://doi.org/10.5281/zenodo.XXXXXXXX}%
%%      {\nolinkurl{doi:10.5281/zenodo.XXXXXXXX}}.

\medskip
\noindent\textbf{Series position.}  Supporting-mathematics repository
of the A=1 DCL series.  Upstream of the planned
\texttt{dcl-formalism} sympy package (which operationalises selected
theorems from this paper as machine-checked primitives) and of the
planned Hilbert-Sixth capstone paper (Paper~I follow-on \#18, which
consumes the discrete-to-continuum limit machinery of
Section~\ref{sec:continuum_limits} for its axiomatisation claim).
Downstream of Paper~I (\emph{Geometry First}, doi:10.5281/zenodo.20078529),
the source-of-record for the bipartite octahedral lattice in its
physical (3+1)-D setting.

\medskip
\noindent\textbf{Status legend (used throughout, including the audit
table in Section~\ref{sec:introduction}).}  \texttt{PASS} -- derivation
complete or experiment confirms the claim to stated precision.
\texttt{PART} -- mechanism demonstrated, full quantitative match
pending; specific gaps documented in the evidence column.
\texttt{STUB} -- analytical result with no full proof yet, or proof
planned but not yet written.  \texttt{FAIL} -- tested and
disconfirmed.

\medskip
\noindent\textbf{Keywords.}  discrete causal lattice; bipartite
octahedral lattice; discrete differential geometry; discrete-to-
continuum limit; Hilbert's sixth problem; foundations of physics.

\newpage

\begin{abstract}
    % abstract.tex
\placeholder{Abstract -- rewrite in finished prose once Phase 2 sections land.}

\noindent
The A=1 Discrete Causal Lattice (DCL) framework establishes the
bipartite octahedral lattice as the substrate of fundamental physics
in (3+1) dimensions (\emph{Geometry First}, Paper~I).  Downstream
work in the DCL series -- the gauge-derivation programme of Paper~II,
the discrete-probability paper closing the gauge-coupling prefactor,
the planned Hilbert-Sixth axiomatisation capstone, and future
kinetic-theory work in the spirit of the Boltzmann / Navier-Stokes
limit -- consumes \emph{more} of the underlying mathematics than
Paper~I exposed, and uses it in arbitrary dimension rather than only
in the (3+1)-D physical setting.  This paper collects the
foundational mathematics those downstream subprojects share, in a
single citable source.

\medskip

\noindent
We define the bipartite octahedral lattice in arbitrary dimension
$n$, characterise its automorphism algebra as a function of $n$
(extending the 3-D 71-dimensional per-site algebra catalogued in
\texttt{dcl-zoo}), develop the discrete differential geometry,
function-space theory, and discrete measure theory the lattice
supports, and -- centrally -- formulate the discrete-to-continuum
limit operators that the Hilbert-Sixth capstone consumes for its
axiomatisation claim.  Two motivating examples appear in
Section~\ref{sec:motivating_examples}: Boltzmann-shaped collision
operators on the lattice, and the lattice-Boltzmann -> Navier-Stokes /
Euler fluid-limit ansatz.  In this first paper these examples are
framing only; full proofs are deferred to subsequent papers in this
repository or in sibling repositories.

\medskip

\noindent
The paper's status across its claims is tracked in
Section~\ref{sec:introduction}, Table~\ref{tab:audit}.  At v0.1-DRAFT,
all toolkit-internal rows are \texttt{STUB}; the artefact at this
version is a paper-shaped outline that locks the conventions and
the sectioning, not a finished body of proofs.

\end{abstract}

\newpage
\tableofcontents
\newpage

%% ============================================================
%% PART I: COMBINATORIAL FOUNDATIONS (THE DIAMOND PROGRESSION)
%% ============================================================
\section{Introduction}
% introduction.tex
% Foundations-toolkit introduction.  Replaces the template's exemplar
% introduction (which demonstrated cross-reference / figure-include
% patterns).  Reference patterns kept; exemplar prose removed.

\label{sec:introduction}

\placeholder{Rewrite once Phase 2 sections land.  This stub fills the
role of an introduction at v0.1: states the problem, the contribution,
and the scope, and surfaces the audit table.}

\subsection{Problem}
\label{subsec:problem}

The A=1 Discrete Causal Lattice (DCL) programme establishes a
bipartite octahedral lattice in (3+1) dimensions as the substrate of
fundamental physics~\cite{example2024}.\footnote{Citations through
this draft refer to placeholder BibTeX entries in
\texttt{paper-bib/references.bib}; real cites land as the prose is
written.}  Paper~I sets out the substrate-priority claim in 3-D
because physics is 3-D.  The mathematics underlying the lattice,
however, is naturally $n$-dimensional, and the series' downstream
work consumes it at that level of generality:

\begin{itemize}
\item Paper~II's gauge-derivation programme uses lattice symmetries
      in a way that does not depend on $n = 3$;
\item the planned Hilbert-Sixth axiomatisation capstone needs
      $n$-dimensional theorems for its universality claim;
\item future kinetic-theory work on the lattice -- the Boltzmann /
      Navier-Stokes programme that Hilbert's 1900 problem singled
      out -- needs discrete-to-continuum limit machinery whose
      natural statement is dimension-agnostic.
\end{itemize}

\noindent
The mathematics is currently scattered across Paper~I's sections,
\texttt{notes/} files in upstream repositories, and the
\texttt{dcl-zoo} catalogue of the 71-dimensional per-site automorphism
algebra.  A single home prevents drift between papers and gives
downstream papers a single citable source.

\subsection{Contribution}
\label{subsec:contribution}

This paper develops the foundations of the A=1 DCL in arbitrary
dimension $n$, organised as follows.

\begin{itemize}
\item \textbf{Combinatorial foundations} (Sections~\ref{sec:lattice_n_dim}
      and~\ref{sec:automorphism_algebra}): the lattice itself in
      dimension $n$, and its automorphism algebra as a function of
      $n$ (recovering the 3-D 71-dimensional algebra at $n = 3$).
\item \textbf{Analytic foundations} (Sections~\ref{sec:discrete_diff_geometry},
      \ref{sec:function_spaces}, and~\ref{sec:discrete_measure_theory}):
      discrete differential geometry, function spaces ($\ell^p$ /
      Sobolev-analogues with $n$-dependent scaling), and discrete
      measure theory.
\item \textbf{The Hilbert-Sixth on-ramp} (Section~\ref{sec:continuum_limits}):
      discrete-to-continuum limit operators -- the most load-bearing
      content for the capstone, where the most mathematical care is
      needed.
\item \textbf{Motivating examples} (Section~\ref{sec:motivating_examples}):
      Boltzmann-shaped collision operators on the lattice, and the
      lattice-Boltzmann -> NS / Euler fluid-limit ansatz.  These
      sections are framing only at v0.1; the full proofs are out of
      scope for this first paper.
\end{itemize}

\subsection{Scope}
\label{subsec:scope}

Out of scope for this first paper:

\begin{itemize}
\item \emph{Physics applications.}  The lattice's role as the
      substrate of fundamental physics is the subject of Papers~I
      and~II and of the discrete-probability paper, not of this
      paper.
\item \emph{Sympy operationalisation.}  Machine-checked versions of
      this paper's theorems will live in the planned
      \texttt{dcl-formalism} package, not here.  This repository is
      paper-only.
\item \emph{Boltzmann / NS proof.}  The lattice kinetic theory and
      the fluid limit appear here as motivating examples
      (Section~\ref{sec:motivating_examples}).  Full proofs of the
      $H$-theorem on the lattice and of the fluid-limit convergence
      are deferred to subsequent papers.
\item \emph{Numerical experiments.}  No code accompanies this
      paper; numerical confirmation of any of its theorems happens
      either in \texttt{dcl-core}-based experiments in sibling repos
      or in the \texttt{dcl-formalism} package's automated
      sympy-verifications.
\end{itemize}

The audit table (Table~\ref{tab:audit}) is the canonical answer to
"has this claim been proved yet?".

%% The audit table is included here so the reader meets it before any
%% claim that depends on a status entry.
% ============================================================
% audit_table.tex
%
% A "system audit" table mapping the paper's claims to the evidence
% backing each one.  This is *the* canonical source of truth for which
% claims are PASS / PART / STUB / FAIL --- prose elsewhere in the paper
% should not contradict it.
%
% Status semantics (also stated in the front-matter of main.tex):
%   PASS  -- derivation complete or experiment confirms the claim to
%            stated precision.
%   PART  -- mechanism demonstrated, full quantitative match pending;
%            specific gaps documented in the evidence column.
%   STUB  -- analytical result with no full proof yet, or planned
%            proof not yet written.
%   FAIL  -- tested and disconfirmed.
%
% Use \texttt{} for code/repo names and file paths inside cells.
% Long cells are fine -- the column widths are tuned for prose.
% ============================================================

\label{tab:audit}

{\scriptsize
\setlength{\tabcolsep}{4pt}
\renewcommand{\arraystretch}{0.95}
\begin{longtable}{@{}p{2.8cm}p{3.8cm}p{3.0cm}p{4.4cm}p{1.0cm}@{}}
\caption{\small System audit: paper claims mapped to supporting
derivations and external references.  Status semantics in the
front-matter (\texttt{PASS} / \texttt{PART} / \texttt{STUB} /
\texttt{FAIL}).  Rows tagged ``inherited'' are upstream-of-record
in Paper~I and act as guardrails the diamond progression must
respect at $d = 3$.  Rows tagged ``motivating example'' describe
scope deferred to subsequent papers.}
\label{tab:audit_caption} \\
\hline
\textbf{Claim} & \textbf{Mechanism} & \textbf{Continuum / formal limit} & \textbf{Evidence} & \textbf{Status} \\
\hline
\endfirsthead

\multicolumn{5}{l}{\textit{(Table~\ref{tab:audit}, continued from previous page)}} \\
\hline
\textbf{Claim} & \textbf{Mechanism} & \textbf{Continuum / formal limit} & \textbf{Evidence} & \textbf{Status} \\
\hline
\endhead

\hline
\multicolumn{5}{r}{\textit{(continued on next page)}} \\
\endfoot

\hline
\endlastfoot

% ── Inherited guardrails from Paper I ──────────────────────────────
Bipartite octahedral lattice at $d = 3$ \emph{(guardrail, Paper~I)}
    & Joint-tick rule on the bipartite octahedral lattice
    & Carries the $(3+1)$-D Lorentz structure
    & Paper~I (\texttt{dcl}), doi:10.5281/zenodo.20078529
    & \texttt{PASS} \\

71-dim per-site automorphism algebra at $d=3$ \emph{(guardrail, \texttt{dcl-zoo})}
    & Per-site automorphisms enumerated; $71$-dim Lie algebra
    & N/A (catalogue)
    & \texttt{dcl-zoo}: generator-zoo repository
    & \texttt{PASS} \\

% ── Diamond progression construction ───────────────────────────────
Diamond progression: d-of-(d+1) simplex vertex construction
    & Combinatorial definition (Definition~\ref{def:diamond_progression}); bipartite by parity
    & Recovers Paper~I's lattice at $d = 3$ (Proposition~\ref{prop:specialisation_paper_i})
    & \S\ref{sec:lattice_n_dim} (paper-text proof pending)
    & \texttt{STUB} \\

% ── Three result-density theorems ──────────────────────────────────
Theorem~1: $\text{coord}(d) = 2d$
    & Count d-of-(d+1) simplex offsets + their negations at each site
    & Distinguishes from chemical-diamond constructions (coord bounded)
    & \S\ref{subsec:theorem_1}, Theorem~\ref{thm:coord_2d} (proof pending)
    & \texttt{STUB} \\

Theorem~2: max Euclidean hop $= \sqrt{d}$
    & Standard simplex-vertex projection into $\mathbb{R}^d$
    & Sets the rescaling parameter for $\epsilon \Lambda_d$ in the continuum limit
    & \S\ref{subsec:theorem_2}, Theorem~\ref{thm:max_hop} (proof pending)
    & \texttt{STUB} \\

Theorem~3: unit-cell volume $= (d+1)^{(d-1)/2}$
    & Simplex-volume formula applied to the offset matrix
    & Enters function-space and measure-theory scaling laws
    & \S\ref{subsec:theorem_3}, Theorem~\ref{thm:unit_cell_volume} (proof pending)
    & \texttt{STUB} \\

% ── Analytic foundations ───────────────────────────────────────────
Automorphism algebra in dimension $d$
    & Generators and commutation relations as functions of $d$; rank closed-form
    & Recovers the 3-D 71-dim per-site algebra at $d=3$
    & \S\ref{sec:automorphism_algebra}; cross-reference \texttt{dcl-zoo}
    & \texttt{STUB} \\

Discrete differential geometry on the diamond progression
    & Discrete exterior calculus: forms, $\mathrm{d}$, codifferential, Stokes; connections / curvature
    & Recovers smooth exterior calculus in the discrete-to-continuum limit (\S\ref{sec:continuum_limits})
    & \S\ref{sec:discrete_diff_geometry} (proof pending)
    & \texttt{STUB} \\

Function spaces on the diamond progression
    & Discrete $L^p$ / $\ell^p$; Sobolev-analogue $\ell^{k,p}$ with $d$-dependent scaling; embedding theorems
    & Standard $L^p$ / $W^{k,p}$ in the limit
    & \S\ref{sec:function_spaces} (proof pending)
    & \texttt{STUB} \\

Discrete measure theory on the diamond progression
    & Counting measure, Haar-analogue invariant under the automorphism algebra, change-of-variables
    & Recovers Lebesgue / Haar measure in the limit
    & \S\ref{sec:discrete_measure_theory} (proof pending)
    & \texttt{STUB} \\

Discrete-to-continuum limit operators \emph{(the Hilbert-Sixth on-ramp)}
    & Lattice spacing $\epsilon$, $\epsilon \to 0$; convergence modes (weak / strong / distributional)
    & Continuum operators in $\mathbb{R}^d$
    & \S\ref{sec:continuum_limits} (proof pending; load-bearing for the planned Hilbert-Sixth capstone)
    & \texttt{STUB} \\

% ── Headline corollary: dimensional selection ──────────────────────
Dimensional selection: $d_{\max} = 1/\delta p_{\min} - 1$
    & Per-site unity constraint at $d+1$ probability slots, each $\geq \delta p_{\min}$
    & Forces $d \leq 1/\delta p_{\min} - 1$
    & \S\ref{sec:dimensional_selection}, Theorem~\ref{thm:dp_min_dmax} (proof pending; depends on Paper~I joint-tick-rule reading)
    & \texttt{STUB} \\

$\delta p_{\min} = 1/4 \Rightarrow d_{\max} = 3$ corollary
    & Direct substitution in Theorem~\ref{thm:dp_min_dmax}
    & N/A (corollary)
    & \S\ref{subsec:dp_min_dmax_3}, Corollary~\ref{cor:dmax_3} (proof immediate once Theorem~\ref{thm:dp_min_dmax} lands)
    & \texttt{STUB} \\

% ── Motivating examples (framing only at v0.2) ─────────────────────
Boltzmann-shaped operators on the diamond progression \emph{(motivating example)}
    & Discrete collision operator; $H$-functional; equilibrium measures
    & Continuum Boltzmann equation in the limit
    & \S\ref{subsec:boltzmann} (framing only; full $H$-theorem proof deferred to a subsequent paper)
    & \texttt{STUB} \\

Fluid (NS / Euler) limit ansatz \emph{(motivating example)}
    & Lattice-Boltzmann $\to$ NS / Euler scaling
    & Navier-Stokes / Euler equations in $\mathbb{R}^d$
    & \S\ref{subsec:fluid_limit} (framing only; full convergence proof deferred)
    & \texttt{STUB} \\

\hline
\end{longtable}
}



\section{The Diamond Progression: d-of-(d+1) Simplex Vertex Construction}
% lattice_n_dim.tex
% Section 2 -- the general bipartite octahedral lattice in dimension n.
% At v0.1: outline-shaped.  Combinatorial definition + parity / two-
% colouring + neighbour structure to be filled in during Phase 2.

\label{sec:lattice_n_dim}

\placeholder{Section body -- write during Phase 2.  Lay out the
combinatorial definition of the bipartite octahedral lattice in
arbitrary dimension $n$; verify that the $n = 3$ specialisation
reproduces Paper~I's lattice.}

\subsection{Definition}
\label{subsec:lattice_definition}

\placeholder{Vertex set, edge set, the bipartite two-colouring, and
the neighbour structure as a function of $n$.  State the lattice as
a graph with site labels in $\mathbb{Z}^n$ (or a suitable index set)
and edges given by a fixed neighbour vector set $\{\pm e_i\}$ plus
the octahedral cross-edges.}

\subsection{Parity / Two-Colouring}
\label{subsec:lattice_parity}

\placeholder{The bipartition into ``even'' and ``odd'' sublattices;
prove that the colouring is well-defined for all $n \geq 1$ and that
edges always cross colour classes.}

\subsection{Specialisation to $n = 3$}
\label{subsec:lattice_specialisation}

\placeholder{Recover Paper~I's 3-D bipartite octahedral lattice as
the $n = 3$ instance of the general construction; cite Paper~I for
the role of this specialisation as the substrate of physics.}

\subsection{Notation conventions}
\label{subsec:lattice_notation}

\placeholder{Standing notation used throughout the rest of the paper:
$\Lambda_n$ for the lattice, $V_n^{+}$ / $V_n^{-}$ for the two colour
classes, $E_n$ for the edge set, $\epsilon$ for the lattice spacing
(used in Section~\ref{sec:continuum_limits} for the limit $\epsilon
\to 0$).}


\section{Three Result-Density Theorems}
% three_theorems.tex
% Section 3 -- the three result-density theorems of the diamond
% progression.  This is the structural spine of the paper.  All STUB
% at v0.2; Phase 2 proves them.

\label{sec:three_theorems}

\placeholder{Section body -- write during Phase~2.  State and prove
the three theorems that characterise the geometry of the diamond
progression as a function of $d$.  Each theorem stands as its own
result; together they distinguish the diamond progression from
nearby constructions.}

\subsection{Theorem 1: Linear Coordination}
\label{subsec:theorem_1}

\begin{theorem}[Linear coordination]
\label{thm:coord_2d}
For every $d \geq 2$, every vertex of $\Lambda_d$ has exactly $2d$
neighbours:
\begin{equation*}
\text{coord}(d) = 2d.
\end{equation*}
\end{theorem}

\begin{proof}
\placeholder{Proof: count the d-of-(d+1) simplex vertex offsets and
their negations.  Each of the $d$ offsets contributes one positive
and one negative edge from any site; total $2d$.  Spell this out
in coordinates once the offset set is fixed.}
\end{proof}

\paragraph{Significance.}
This is the most consequential of the three theorems.  Chemical-
diamond and face-centred-cubic-style constructions have coordination
that stays bounded (typically~$4$ or~$12$) as $d$ grows; the diamond
progression's linear coordination is the structural feature that
distinguishes it from those constructions and that underwrites the
dimensional-selection corollary in
Section~\ref{sec:dimensional_selection}.

\subsection{Theorem 2: Maximum Euclidean Hop}
\label{subsec:theorem_2}

\begin{theorem}[Maximum hop length]
\label{thm:max_hop}
For every $d \geq 2$, the longest single edge (in Euclidean distance,
under the standard projection of the regular $(d+1)$-simplex into
$\mathbb{R}^d$) has length $\sqrt{d}$.
\end{theorem}

\begin{proof}
\placeholder{Proof: the longest edge corresponds to the simplex
vertex furthest from a chosen base vertex.  Compute the Euclidean
distance in the standard projection; it works out to $\sqrt{d}$.
Spell out the projection conventions and the explicit distance
computation.}
\end{proof}

\paragraph{Significance.}
Combined with Theorem~1, this implies that as $d$ grows the lattice
becomes \emph{denser in hops per unit Euclidean distance} -- there
are $2d$ edges per vertex but only one of length $\sqrt{d}$; the rest
are shorter.  Propagation through the $d = 3$ lattice is therefore
slower (more hops per unit physical distance) than through higher-$d$
instances.  This observation is the seed of the \emph{essay}
companion ``the progression is fractal'' tracked in `raw_thoughts.md`.

\subsection{Theorem 3: Unit-Cell Volume}
\label{subsec:theorem_3}

\begin{theorem}[Unit-cell volume]
\label{thm:unit_cell_volume}
For every $d \geq 2$, the volume of the unit cell of $\Lambda_d$ is
\begin{equation*}
\text{vol}(d) = (d+1)^{(d-1)/2}.
\end{equation*}
\end{theorem}

\begin{proof}
\placeholder{Proof: standard simplex-volume formula applied to the
$(d+1)$-vertex projection in $\mathbb{R}^d$.  $(d+1)^{(d-1)/2}$
follows from the determinant of the offset matrix.  Spell out the
matrix and the determinant computation.}
\end{proof}

\paragraph{Significance.}
Unit-cell volume grows super-exponentially in $d$ -- $\text{vol}(3)
= 4 \cdot \sqrt{2} \approx 5.66$, $\text{vol}(4) = 25$,
$\text{vol}(5) = 6^2 \sqrt{6} \approx 88$.  Combined with Theorems~1
and~2, this gives the density / scaling profile that the discrete-to-
continuum limit machinery (Section~\ref{sec:continuum_limits})
must respect.

\subsection{Why These Three}
\label{subsec:why_three_theorems}

\placeholder{The three theorems characterise the diamond progression
at three distinct geometric levels: combinatorial
(Theorem~\ref{thm:coord_2d}), Euclidean
(Theorem~\ref{thm:max_hop}), and volumetric
(Theorem~\ref{thm:unit_cell_volume}).  A fourth theorem of the same
calibre would be one of: the explicit shape of the unit cell beyond
volume (boundary structure); the second-eigenvalue / spectral gap
of the adjacency operator; the diameter / girth of $\Lambda_d$
restricted to a finite torus.  Any of these can be added in a
follow-on paper without disturbing the three-theorem spine of this
one.}

\subsection{Implications for the Rest of the Paper}
\label{subsec:three_theorem_implications}

\begin{itemize}
\item Theorem~\ref{thm:coord_2d} ($\text{coord} = 2d$) is the
      key input to the dimensional-selection corollary
      (Section~\ref{sec:dimensional_selection}).
\item Theorem~\ref{thm:max_hop} (max hop $= \sqrt{d}$) sets the
      scale parameter for the rescaled lattice $\epsilon \Lambda_d$
      in the discrete-to-continuum limit
      (Section~\ref{sec:continuum_limits}).
\item Theorem~\ref{thm:unit_cell_volume} (vol $= (d+1)^{(d-1)/2}$)
      enters the function-space scaling laws
      (Section~\ref{sec:function_spaces}) and the discrete-measure
      definitions
      (Section~\ref{sec:discrete_measure_theory}).
\end{itemize}


\section{Automorphism Algebra in Dimension $d$}
% automorphism_algebra.tex
% Section 3 -- automorphism algebra of the bipartite octahedral lattice
% in dimension n.  Extends the 3-D 71-dimensional per-site algebra
% catalogued in dcl-zoo; gives generators and commutation relations as
% a function of n, with rank closed-form.

\label{sec:automorphism_algebra}

\placeholder{Section body -- write during Phase 2.  Develop the
per-site automorphism algebra as a function of $n$; identify the
generators (rotations, reflections, parity-preserving / parity-
flipping pieces); state and prove commutation relations.}

\subsection{Generators}
\label{subsec:autom_generators}

\placeholder{Enumerate the per-site automorphism generators in
arbitrary dimension.  The octahedral / hyperoctahedral structure
gives a natural set; characterise it intrinsically rather than by
appeal to the $n = 3$ enumeration.}

\subsection{Commutation Relations}
\label{subsec:autom_commutators}

\placeholder{State the commutation relations among the generators;
identify the structure constants as functions of $n$.}

\subsection{Rank and Dimension}
\label{subsec:autom_rank}

\placeholder{Closed-form expression for the dimension of the per-site
algebra as a function of $n$.  Verify that $n = 3$ recovers the
71-dimensional algebra catalogued in \texttt{dcl-zoo}.}

\subsection{Specialisation to $n = 3$}
\label{subsec:autom_specialisation}

\placeholder{Recover the 3-D 71-dimensional per-site algebra as
the $n = 3$ instance.  Cross-reference \texttt{dcl-zoo} for the
explicit generator catalogue at $n = 3$.}


%% ============================================================
%% PART II: ANALYTIC FOUNDATIONS
%% ============================================================
\section{Discrete Differential Geometry on the Diamond Progression}
% discrete_diff_geometry.tex
% Section 4 -- discrete differential geometry on the bipartite
% octahedral lattice.  Discrete exterior calculus (forms, d, codifferential),
% discrete connections and curvature, integration by parts.

\label{sec:discrete_diff_geometry}

\placeholder{Section body -- write during Phase 3.  Develop the
discrete exterior calculus on the lattice: discrete forms, discrete
exterior derivative $d$, discrete codifferential, Stokes' theorem
analogue, discrete connections, and discrete curvature.}

\subsection{Discrete Forms}
\label{subsec:dec_forms}

\placeholder{Define discrete $k$-forms on the lattice as functions on
$k$-cells (vertices, edges, plaquettes, \ldots).  Address the
bipartite colour structure: how forms transform under colour-class
exchange.}

\subsection{Discrete Exterior Derivative}
\label{subsec:dec_d}

\placeholder{Define the discrete exterior derivative $d \colon
\Omega^k \to \Omega^{k+1}$ on the lattice.  Prove $d \circ d = 0$.}

\subsection{Discrete Codifferential and Inner Product}
\label{subsec:dec_codifferential}

\placeholder{Inner product on discrete forms (using the discrete
measure of Section~\ref{sec:discrete_measure_theory}).  Define the
codifferential $d^\ast$ as the formal adjoint of $d$.  State the
discrete Hodge decomposition.}

\subsection{Discrete Stokes' Theorem and Integration by Parts}
\label{subsec:dec_stokes}

\placeholder{Lattice-Stokes identity and the integration-by-parts
formula that follows.  This is the load-bearing identity for the
function-space theory in Section~\ref{sec:function_spaces}.}

\subsection{Discrete Connections and Curvature}
\label{subsec:dec_connections}

\placeholder{Discrete connections on the lattice (gauge / parallel-
transport viewpoint); discrete curvature as the failure of the
plaquette holonomy to commute.  Cross-reference Paper~II's gauge
derivation as a downstream consumer.}


\section{Function Spaces on the Diamond Progression}
% function_spaces.tex
% Section 5 -- function spaces on the bipartite octahedral lattice.
% Discrete L^p / ell^p, Sobolev-analogue spaces ell^{k,p}, n-dependent
% scaling, embedding theorems.

\label{sec:function_spaces}

\placeholder{Section body -- write during Phase 3.  Define the
discrete $\ell^p$ and Sobolev-analogue $\ell^{k,p}$ spaces on the
lattice; state and prove $n$-dependent scaling, density, and
embedding theorems.}

\subsection{Discrete $\ell^p$ Spaces}
\label{subsec:lp_spaces}

\placeholder{Define $\ell^p(\Lambda_n)$ for $1 \leq p \leq \infty$
using the counting measure on $\Lambda_n$.  Standard
Banach-space properties: completeness, separability for $p <
\infty$, duality.}

\subsection{Sobolev-Analogue Spaces}
\label{subsec:sobolev_analogue}

\placeholder{Define $\ell^{k,p}(\Lambda_n)$ as the space of lattice
functions whose discrete-derivatives up to order $k$ (via the
discrete exterior derivative of
Section~\ref{sec:discrete_diff_geometry}) are in $\ell^p$.  State
the basic structural theorems (completeness, density of
finitely-supported functions, duality).}

\subsection{Scaling and Embedding Theorems}
\label{subsec:embeddings}

\placeholder{$n$-dependent scaling of the norms under
lattice-spacing rescaling $\epsilon$.  State discrete analogues of
the Sobolev embedding $W^{k,p} \hookrightarrow L^q$ with explicit
$n$-dependent exponent relations.  These are the embeddings that
will pass to the continuum limit in
Section~\ref{sec:continuum_limits}.}

\subsection{Compact Embeddings}
\label{subsec:compact_embeddings}

\placeholder{Compactness statements analogous to Rellich-Kondrachov
in the discrete setting; the role of the lattice-spacing parameter
$\epsilon$ in controlling compactness.}


\section{Discrete Measure Theory on the Diamond Progression}
% discrete_measure_theory.tex
% Section 6 -- discrete measure theory on the bipartite octahedral
% lattice.  Counting measures, Haar-analogue measures invariant under
% the automorphism algebra, change-of-variables, discrete Radon-Nikodym.

\label{sec:discrete_measure_theory}

\placeholder{Section body -- write during Phase 3.  Develop the
measure theory needed to underwrite the function-space theory in
Section~\ref{sec:function_spaces} and the integration-by-parts
identity in Section~\ref{sec:discrete_diff_geometry}.}

\subsection{Counting Measure}
\label{subsec:counting_measure}

\placeholder{The basic counting measure $\mu_\#$ on $\Lambda_n$;
its restriction to colour classes $V_n^{\pm}$; sigma-algebra
generated by finite sets.}

\subsection{Haar-Analogue Measures}
\label{subsec:haar_analogue}

\placeholder{Measures invariant under the per-site automorphism
algebra of Section~\ref{sec:automorphism_algebra}; uniqueness up to
scaling; the analogue of Haar's theorem for the lattice.}

\subsection{Change-of-Variables}
\label{subsec:change_of_variables}

\placeholder{Change-of-variables formula for measure-preserving
maps and for the lattice-spacing rescaling $\epsilon \mapsto
\lambda \epsilon$.  Underwrites the limit-rescaling in
Section~\ref{sec:continuum_limits}.}

\subsection{Discrete Radon-Nikodym}
\label{subsec:radon_nikodym}

\placeholder{Discrete analogue of the Radon-Nikodym theorem: when a
measure $\nu$ on $\Lambda_n$ has a density with respect to
$\mu_\#$, and the explicit formula for the density.}


%% ============================================================
%% PART III: THE HILBERT-SIXTH ON-RAMP
%% ============================================================
\section{Discrete-to-Continuum Limit Operators}
% continuum_limits.tex
% Section 8 -- discrete-to-continuum limit operators.  The Hilbert-Sixth
% on-ramp.  Lattice-spacing parameter epsilon, epsilon -> 0 limit,
% convergence modes, consistency conditions tying discrete to
% continuum operators.

\label{sec:continuum_limits}

\placeholder{Section body -- write during Phase 4.  This is the
load-bearing section for the planned Hilbert-Sixth capstone.  Develop
the discrete-to-continuum limit machinery: parameter, topology,
convergence modes, consistency theorems.  The most mathematical
care in the paper goes here.}

\subsection{The Limit Parameter}
\label{subsec:limit_parameter}

\placeholder{Introduce the lattice-spacing parameter $\epsilon > 0$;
the rescaled lattice $\epsilon \Lambda_d \subset \mathbb{R}^d$; the
$\epsilon \to 0$ limit as the discrete-to-continuum passage.}

\subsection{Convergence Modes}
\label{subsec:convergence_modes}

\placeholder{Define what it means for a sequence of discrete
operators $T_\epsilon \colon \ell^p(\epsilon \Lambda_d) \to
\ell^p(\epsilon \Lambda_d)$ to converge to a continuum operator
$T \colon L^p(\mathbb{R}^d) \to L^p(\mathbb{R}^d)$.  Three modes:
strong, weak, and distributional.  Relations among them.}

\subsection{Consistency Conditions}
\label{subsec:consistency_conditions}

\placeholder{Conditions tying discrete operators to their continuum
limits.  The discrete exterior derivative $\mathrm{d}$
(Section~\ref{sec:discrete_diff_geometry}) converges to the
continuum $\mathrm{d}$ in $L^p$ for appropriate function classes;
the discrete Hodge Laplacian $\Delta_{\Lambda_d}$ converges to the
continuum Laplacian; \ldots.  These are the consistency theorems
the Hilbert-Sixth capstone will cite.}

\subsection{Function-Space Convergence}
\label{subsec:function_space_convergence}

\placeholder{Convergence of the function spaces themselves:
$\ell^p(\epsilon \Lambda_d) \to L^p(\mathbb{R}^d)$ in an
appropriate sense; Sobolev-embedding constants stabilising in the
limit.}

\subsection{Why This is the Hilbert-Sixth On-Ramp}
\label{subsec:hilbert_sixth_onramp}

\placeholder{Hilbert's 1900 problem specifically singled out
Boltzmann's kinetic theory and the rigorous derivation of fluid
equations as the centerpiece of an axiomatised physics.  Both
require a discrete-to-continuum limit treated with full rigour.
This section provides the apparatus the capstone synthesises into
the axiomatisation claim.  Cross-reference: the motivating examples
in Section~\ref{sec:motivating_examples} show what the apparatus is
used \emph{for}; the dimensional-selection corollary in
Section~\ref{sec:dimensional_selection} is the other principal
consumer.}


%% ============================================================
%% PART IV: DIMENSIONAL SELECTION (HEADLINE COROLLARY)
%% ============================================================
\section{Dimensional Selection: $d_{\max} = 1/\delta p_{\min} - 1$}
% dimensional_selection.tex
% Section 9 (NEW in v0.2) -- the headline corollary: a minimum
% probability quantum δp_min forces d_max = 1/δp_min - 1.  At δp_min
% = 1/4, d_max = 3.  This is the candidate mechanism that selects
% the universe's dimension rather than permitting it.

\label{sec:dimensional_selection}

\placeholder{Section body -- write during Phase~2.  State and prove
the dimensional-selection corollary; spell out the
$\delta p_\min = 1/4 \Rightarrow d_\max = 3$ specialisation; document
the cross-reference to \texttt{dcl-delta-p-min} for the empirical
investigation of $\delta p_\min$'s value.}

\subsection{Setup: A Minimum Probability Quantum}
\label{subsec:dp_min_setup}

\placeholder{Take $\delta p_\min > 0$ as a hypothesis: the A=1
framework admits no probability per site smaller than $\delta p_\min$.
The hypothesis can be realised in (at least) two operationally
distinct ways:

\begin{itemize}
\item A \emph{clamp} on continuous amplitudes: $|psi(x)|^2 \geq
   \delta p_\min$ when non-zero.  Implementation: the
   \texttt{prob\_floor} parameter on \texttt{dcl\_core.core}
   (planned \texttt{dcl-core v0.2.0}).
\item An \emph{intrinsic granularity}: probability is represented as
   a multiple of $\delta p_\min$.  Implementation: the integer-token
   engine \texttt{dcl\_core.core3d} with token budget
   $N = 1/\delta p_\min$.
\end{itemize}

The cross-engine equivalence question -- whether these two
realisations produce the same physics -- is the open question of
Phase~1 in \texttt{dcl-delta-p-min}.  This section's result is
operationally engine-agnostic: it depends only on the *existence* of
a $\delta p_\min > 0$, not on its realisation.}

\subsection{The Dimensional-Selection Corollary}
\label{subsec:dp_min_dmax_corollary}

\begin{theorem}[Dimensional selection from $\delta p_\min$]
\label{thm:dp_min_dmax}
If a minimum probability quantum $\delta p_\min > 0$ exists in the
A=1 framework, then the diamond progression $\Lambda_d$ can sustain
the $\mathcal{A} = 1$ unity axiom only for
\begin{equation*}
d \leq d_{\max} = \frac{1}{\delta p_\min} - 1.
\end{equation*}
\end{theorem}

\begin{proof}
\placeholder{Proof: at any site $x$ of $\Lambda_d$, the unity
constraint $\mathcal{A} = 1$ requires the per-site probability mass
to distribute over $x$ together with all the configurations
participating in the joint-tick rule.  By
Theorem~\ref{thm:coord_2d}, the site has $2d$ neighbours; combined
with the joint-tick-rule's structure, $d + 1$ distinct probability
slots receive non-zero mass.  Each slot must carry probability at
least $\delta p_\min$, so
$(d+1) \cdot \delta p_\min \leq 1$, giving
$d \leq 1/\delta p_\min - 1$.  Spell out the joint-tick-rule
participation count once the construction is fixed in
Section~\ref{sec:lattice_n_dim}.}
\end{proof}

\paragraph{Note on the proof structure.}
\todomark{motivating-example: the exact slot-count ``$d+1$'' depends
on a specific reading of the joint-tick rule that pulls from Paper~I
plus the diamond-progression structure of
Section~\ref{sec:lattice_n_dim}.  Phase~2 work must verify that
this count is precise rather than a sketch; alternative readings
(e.g.\ $2d+1$ slots) give different $d_\max$ formulae.  The
formula $d_\max = 1/\delta p_\min - 1$ is the user's stated result
from `raw\_thoughts.md` 2026-05-21 \texttt{[result-2]}; this paper
re-derives and confirms it.}

\subsection{Specialisation: $\delta p_\min = 1/4 \Rightarrow d_{\max} = 3$}
\label{subsec:dp_min_dmax_3}

\begin{corollary}
\label{cor:dmax_3}
If $\delta p_\min = 1/4$, then $d_{\max} = 3$.
\end{corollary}

\begin{proof}
Direct substitution into Theorem~\ref{thm:dp_min_dmax}:
$d_{\max} = 1/(1/4) - 1 = 4 - 1 = 3$.
\end{proof}

\paragraph{Significance.}
This corollary is the candidate mechanism that \emph{forces} the
universe to be three-dimensional rather than permitting it.  If the
A=1 framework forces $\delta p_\min = 1/4$ (via a token-conservation
constraint or analogous structural argument), then it forces $d =
3$.  The framework is no longer compatible with our observed
universe by accident -- it is compatible because the unique
admissible $d$ is $3$.

\paragraph{Connection to Paper~I.}
Paper~I established the bipartite octahedral lattice as the
substrate of physics in $(3+1)$ dimensions.  At the time of
Paper~I, the specific choice of $d = 3$ was \emph{motivated} by
observation -- physics is 3D, so the lattice is 3D.  This
corollary inverts the relationship: given $\delta p_\min$, the
framework forces $d = 3$.  Observational compatibility becomes
a structural consequence rather than a motivating constraint.

\subsection{The Empirical Question: Is $\delta p_\min = 1/4$?}
\label{subsec:dp_min_value_question}

\placeholder{The corollary above is mathematical: it derives
$d_\max$ from $\delta p_\min$ given the diamond progression's
structure.  Whether the framework forces $\delta p_\min$ to take a
specific value (and which value) is a separate question -- empirical
and / or structural-from-token-conservation.

The companion repository \texttt{dcl-delta-p-min} runs the
investigation.  Briefly: a 4-cell numerical grid (two engines
crossed with two $\delta p_\min$ regimes) measures how the
Arnold-tongue width and hydrogen Bohr-radius recovery depend on
$\delta p_\min$; the largest $\delta p_\min$ at which Paper~I's
PASS audit rows remain recoverable is a framework-internal upper
bound.  The planned discrete-probability paper plausibly pins
$\delta p_\min$'s natural value via a token-conservation argument
(Outcome~D of \texttt{dcl-delta-p-min}'s Phase~4 stance decision).
At $\delta p_\min = 1/4$, Corollary~\ref{cor:dmax_3} forces $d_\max
= 3$ exactly; at other values, the dimensional ceiling shifts.

For this paper, the result stands as a \emph{conditional theorem}:
if $\delta p_\min$ exists and equals $1/4$, then $d_\max = 3$.
The conditional is non-empty -- the diamond progression admits
$\delta p_\min$ structurally, and the value $1/4$ is the empirical
target of \texttt{dcl-delta-p-min}.}

\subsection{Implications}
\label{subsec:dimensional_selection_implications}

\begin{itemize}
\item \emph{The Hilbert-Sixth capstone}
      (\texttt{dcl-paper-XX-hilbert-sixth}, planned) consumes this
      corollary as part of its universality argument: the
      axiomatisation forces $d = 3$.
\item \emph{The planned discrete-probability paper} fixes
      $\delta p_\min$'s value (Outcome~D of the
      \texttt{dcl-delta-p-min} investigation); this paper's
      corollary then becomes a parameter-free statement.
\item \emph{The aspirational F4 sequel} (`A=1 forces linear
      coordination', tracked in
      \texttt{physics-research/notes/raw\_thoughts.md}) asks
      whether the diamond progression is the \emph{only} bipartite
      lattice family admitting A=1 non-degenerately.  If F4 lands,
      Corollary~\ref{cor:dmax_3} sharpens further: not just
      ``$d_\max = 3$ for the diamond progression'' but ``$d_\max
      = 3$ for any lattice family admitting A=1''.
\end{itemize}


%% ============================================================
%% PART V: MOTIVATING EXAMPLES (FRAMING ONLY)
%% ============================================================
\section{Motivating Examples: Boltzmann and the Fluid Limit}
% motivating_examples.tex
% Section 10 -- motivating examples (framing only).  Boltzmann-shaped
% collision operators on the lattice and the lattice-Boltzmann ->
% NS / Euler fluid-limit ansatz.  Full proofs are out of scope for
% this paper; deferred to subsequent papers in this repo or in sibling
% repos.

\label{sec:motivating_examples}

\placeholder{Section body -- write during Phase 5.  This section is
\emph{framing only}: it states what the foundations toolkit
developed in Sections~\ref{sec:lattice_n_dim}--\ref{sec:continuum_limits}
will eventually be used for.  Full proofs of the H-theorem on
$\Lambda_d$ and of the fluid-limit convergence are deferred to
subsequent papers.}

\subsection{Boltzmann-Shaped Operators on the Diamond Progression}
\label{subsec:boltzmann}

\placeholder{Define a Boltzmann-shaped collision operator on
$\Lambda_d$ (or its rescaling $\epsilon \Lambda_d$): a binary /
$k$-ary collision rule that respects the bipartite structure and
the automorphism algebra of
Section~\ref{sec:automorphism_algebra}.  Define the lattice
$H$-functional.  State (without proving) the expected $H$-theorem
and the expected equilibrium measures.  Cross-reference to a future
paper for the full proof.}

\paragraph{Motivating-example status.}
\todomark{motivating-example: the $H$-theorem on $\Lambda_d$ is
asserted here but not proved.  Full proof scheduled for a
subsequent paper.  See Table~\ref{tab:audit}, ``Boltzmann-shaped
operators'' row, which is by design \texttt{STUB}.}

\subsection{Fluid (Navier-Stokes / Euler) Limit Ansatz}
\label{subsec:fluid_limit}

\placeholder{State the lattice-Boltzmann $\to$ NS / Euler limit
ansatz: the scaling of $\epsilon$ and of the collision-rule
parameters that should recover the continuum fluid equations in
the limit (via the convergence machinery of
Section~\ref{sec:continuum_limits}).  Identify the technical
ingredients the convergence proof would need.  Sketch the
expected statement.  Cross-reference to a future paper for the
full proof.}

\paragraph{Motivating-example status.}
\todomark{motivating-example: the lattice-Boltzmann $\to$ NS / Euler
convergence is asserted here but not proved.  Full proof scheduled
for a subsequent paper.  See Table~\ref{tab:audit}, ``Fluid
(NS / Euler) limit ansatz'' row, which is by design \texttt{STUB}.}

\subsection{Why These Examples}
\label{subsec:why_these_examples}

\placeholder{Hilbert's 1900 problem singled out Boltzmann's kinetic
theory and the rigorous Boltzmann $\to$ NS / Euler limit as the
centerpiece of an axiomatised physics.  Treating them as motivating
examples here -- rather than as proof targets of this paper --
keeps this first paper focused on the foundations toolkit while
making the eventual use of that toolkit visible.  The full proofs
of the $H$-theorem on $\Lambda_d$ and of the fluid-limit
convergence are scheduled for subsequent papers in this repository
or in sibling repositories.}


\section{Conclusion}
\input{sections/conclusion}

\section*{Acknowledgements}
\input{sections/acknowledgements}

%% ============================================================
%% APPENDICES
%% ============================================================
\appendix

\section{Symbols}
% symbols.tex
% Appendix -- Symbols and meanings.  Reference list of mathematical
% symbols introduced in this paper.  Update this list whenever a new
% standing symbol enters the body of the paper.
%
% Section column is the section where the symbol is first used in its
% standing meaning.  Many symbols are reused in later sections; the
% Section column points to the *definition*, not the most recent use.

\label{sec:symbols}

\placeholder{Update entries whenever a new standing symbol enters the
body of the paper.  Many entries below are placeholders pinned to
their planned section; the meaning is fixed by the
\texttt{scope-and-outline} note (\texttt{notes/dcl\_mathematics\_scope\_and\_outline.md}).}

{\small
\setlength{\tabcolsep}{6pt}
\renewcommand{\arraystretch}{1.10}
\begin{longtable}{@{}p{2.6cm}p{8.4cm}p{3.0cm}@{}}
\caption{\small Symbols and their meanings, with the section where
each symbol is first defined.  Symbols introduced only inside a
single proof or example (e.g.\ dummy summation indices) are not
listed.}
\label{tab:symbols} \\
\hline
\textbf{Symbol} & \textbf{Meaning} & \textbf{Section} \\
\hline
\endfirsthead

\multicolumn{3}{l}{\textit{(Table~\ref{tab:symbols}, continued from previous page)}} \\
\hline
\textbf{Symbol} & \textbf{Meaning} & \textbf{Section} \\
\hline
\endhead

\hline
\multicolumn{3}{r}{\textit{(continued on next page)}} \\
\endfoot

\hline
\endlastfoot

% ── Lattice and ambient spaces ─────────────────────────────────────
\multicolumn{3}{@{}l}{\textit{Lattice and ambient spaces}} \\
$n$
    & Lattice dimension; free parameter throughout the paper.  Specialised to $n = 3$ only when matching Paper~I's physical setting.
    & \S\ref{sec:lattice_n_dim} \\
$\Lambda_n$
    & The bipartite octahedral lattice in dimension $n$.
    & \S\ref{sec:lattice_n_dim} \\
$V_n$
    & Vertex set of $\Lambda_n$, identified with $\mathbb{Z}^n$ (provisional).
    & \S\ref{subsec:lattice_definition} \\
$V_n^{+}, V_n^{-}$
    & The two colour classes of the bipartition; $V_n = V_n^{+} \sqcup V_n^{-}$.
    & \S\ref{subsec:lattice_parity} \\
$E_n$
    & Edge set of $\Lambda_n$: axial $\pm e_i$ edges plus octahedral cross-edges.
    & \S\ref{subsec:lattice_definition} \\
$\epsilon$
    & Lattice-spacing parameter ($\epsilon > 0$); $\epsilon \to 0$ gives the discrete-to-continuum limit.
    & \S\ref{subsec:limit_parameter} \\
$\epsilon \Lambda_n$
    & Rescaled lattice $\subset \mathbb{R}^n$ with spacing $\epsilon$.
    & \S\ref{subsec:limit_parameter} \\
$\mathbb{R}^n$, $\mathbb{Z}^n$
    & Continuum ambient space; integer lattice (used as vertex-set model for $\Lambda_n$).
    & \S\ref{subsec:lattice_definition} \\

% ── Symmetries ─────────────────────────────────────────────────────
\multicolumn{3}{@{}l}{\textit{Symmetries}} \\
$B_n$
    & Hyperoctahedral group: signed permutations of the $n$ coordinate axes; acts on $\Lambda_n$.
    & \S\ref{subsec:autom_generators} \\
$\mathfrak{a}_n$
    & The per-site automorphism algebra of $\Lambda_n$; provisional notation. At $n = 3$, $\dim \mathfrak{a}_3 = 71$ (catalogued in \texttt{dcl-zoo}).
    & \S\ref{sec:automorphism_algebra} \\

% ── Discrete differential geometry ─────────────────────────────────
\multicolumn{3}{@{}l}{\textit{Discrete differential geometry}} \\
$\Omega^k(\Lambda_n)$
    & Discrete $k$-forms on $\Lambda_n$: functions on oriented $k$-cells.
    & \S\ref{subsec:dec_forms} \\
$d$
    & Discrete exterior derivative $\Omega^k \to \Omega^{k+1}$.  Satisfies $d \circ d = 0$.
    & \S\ref{subsec:dec_d} \\
$d^\ast$
    & Discrete codifferential: formal adjoint of $d$ with respect to the inner product on forms.
    & \S\ref{subsec:dec_codifferential} \\
$\langle \cdot, \cdot \rangle$
    & Inner product on $\Omega^k(\Lambda_n)$ defined via the discrete measure.
    & \S\ref{subsec:dec_codifferential} \\
$\Delta = -d^\ast d - d d^\ast$
    & Discrete Hodge Laplacian on $\Omega^k(\Lambda_n)$.  Converges to the continuum Laplacian in the $\epsilon \to 0$ limit.
    & \S\ref{subsec:consistency_conditions} \\

% ── Function spaces ────────────────────────────────────────────────
\multicolumn{3}{@{}l}{\textit{Function spaces}} \\
$\ell^p(\Lambda_n)$
    & Discrete Lebesgue space, $1 \leq p \leq \infty$.
    & \S\ref{subsec:lp_spaces} \\
$\ell^{k,p}(\Lambda_n)$
    & Discrete Sobolev-analogue space: lattice functions with discrete derivatives (via $d$) up to order $k$ in $\ell^p$.
    & \S\ref{subsec:sobolev_analogue} \\
$\|f\|_{p, \epsilon}$
    & $\epsilon$-rescaled $\ell^p$-norm: $\|f\|_{p, \epsilon}^p = \epsilon^n \sum_{x} |f(x)|^p$.  Converges to $\|f\|_{L^p(\mathbb{R}^n)}$ in the $\epsilon \to 0$ limit.
    & \S\ref{subsec:embeddings} \\
$L^p(\mathbb{R}^n)$, $W^{k,p}(\mathbb{R}^n)$
    & Continuum Lebesgue and Sobolev spaces; targets of the discrete-to-continuum limit.
    & \S\ref{subsec:function_space_convergence} \\

% ── Measure theory ─────────────────────────────────────────────────
\multicolumn{3}{@{}l}{\textit{Measure theory}} \\
$\mu_\#$
    & Counting measure on $\Lambda_n$: $\mu_\#(A) = |A|$ for finite $A \subset \Lambda_n$.
    & \S\ref{subsec:counting_measure} \\
$\mu_{\#, \epsilon}$
    & $\epsilon$-rescaled counting measure on $\epsilon \Lambda_n$: $\mu_{\#, \epsilon} = \epsilon^n \mu_\#$.  Converges to Lebesgue measure on $\mathbb{R}^n$ in the limit.
    & \S\ref{subsec:counting_measure} \\

% ── Kinetic theory (motivating examples) ───────────────────────────
\multicolumn{3}{@{}l}{\textit{Kinetic theory (motivating examples, Section~\ref{sec:motivating_examples})}} \\
$V$
    & Discrete velocity space attached to each lattice site.  Provisionally, the neighbour-vector set of $\Lambda_n$.
    & \S\ref{subsec:boltzmann} \\
$f(x, v, t)$
    & Lattice kinetic distribution function: density of particles at $x \in \Lambda_n$ with velocity $v \in V$ at time $t$.
    & \S\ref{subsec:boltzmann} \\
$f_{\text{eq}}$
    & Local equilibrium (Maxwellian-analogue) distribution; parametrised by the conserved moments.
    & \S\ref{subsec:boltzmann} \\
$Q[f]$
    & Boltzmann-shaped collision operator on the lattice.
    & \S\ref{subsec:boltzmann} \\
$H[f]$
    & Lattice $H$-functional: $H[f] = \sum_{x, v} f(x, v) \log f(x, v)$.  Conjectured non-increasing under the lattice kinetic dynamics.
    & \S\ref{subsec:boltzmann} \\
$\rho$, $u$, $e$
    & Fluid density, velocity, energy (the conserved moments of $f$); arguments of $f_{\text{eq}}$.
    & \S\ref{subsec:fluid_limit} \\

\hline
\end{longtable}
}


\section{Terms}
% terms.tex
% Appendix -- Terms and meanings.  Glossary of technical / framework
% terminology used in this paper.  Alphabetical within categories;
% categories appear in a fixed order (framework -> mathematics ->
% kinetic theory).  Cross-references to upstream papers and external
% references are kept short; the bibliography carries the full cites.

\label{sec:terms}

\placeholder{Update entries whenever a term enters or shifts its
standing meaning in the body of the paper.  Definitions below match
the working notes in \texttt{notes/} -- when the two diverge, the
notes are the working source and this appendix gets the synchronised
phrasing at edit time.}

\subsection*{Framework terms}

\begin{description}
\item[A=1 (Discrete Causal Lattice)] The A=1 DCL is the framework
   established in Paper~I (\emph{Geometry First}).  Physics emerges
   from three axioms applied to a bipartite octahedral lattice:
   \emph{Unity} (amplitude is conserved at every tick, hence A=1),
   \emph{Locality} (information propagates only between adjacent
   nodes), and \emph{Phase} ($U(1)$ on a per-particle internal
   clock).  See Paper~I for the full statement.

\item[Audit table] A table mapping the paper's claims to the
   evidence backing each one.  Each row has a status in
   $\{\texttt{PASS}, \texttt{PART}, \texttt{STUB}, \texttt{FAIL}\}$.
   The table is canonical: prose elsewhere in the paper does not
   contradict it.  In this paper, the audit table appears in
   \S\ref{sec:introduction}, Table~\ref{tab:audit}.

\item[Bipartite octahedral lattice] The lattice $\Lambda_n$ defined
   combinatorially in \S\ref{sec:lattice_n_dim}.  In dimension
   $n = 3$, recovers Paper~I's substrate lattice.  Bipartite: vertex
   set splits into two colour classes $V_n^{+}, V_n^{-}$.  Octahedral:
   the neighbour structure at each site is the octahedral
   coordination figure (or its $n$-D generalisation).

\item[\texttt{dcl-formalism} (package)] Planned sympy-extending
   package operationalising selected theorems of this paper as
   machine-checked primitives.  Distinct from this repository:
   \texttt{dcl-mathematics} proves the theorems; \texttt{dcl-formalism}
   implements them.

\item[\texttt{dcl-zoo} (generator zoo)] Catalogue repository
   enumerating the 71-dimensional per-site automorphism algebra at
   $n = 3$.  See \S\ref{subsec:autom_specialisation}.

\item[Hilbert's Sixth Problem] Problem 6 from Hilbert's 1900 list:
   \emph{``treat in the same manner, by means of axioms, those
   physical sciences in which already today mathematics plays an
   important part.''}  Singled out Boltzmann's kinetic theory and
   the rigorous derivation of fluid equations as the centerpiece.
   The planned series capstone (\texttt{dcl-paper-XX-hilbert-sixth},
   Paper~I follow-on \#18) consumes this paper's
   discrete-to-continuum limit machinery
   (\S\ref{sec:continuum_limits}) for its axiomatisation claim.

\item[Hilbert-Sixth on-ramp] Informal label for
   \S\ref{sec:continuum_limits} of this paper: the
   discrete-to-continuum limit machinery whose rigorous statement
   the Hilbert-Sixth capstone synthesises into the axiomatisation
   claim.

\item[Joint-tick rule] Paper~I's rule for advancing the bipartite
   lattice's two colour classes in synchrony at every tick.  Used
   in \S\ref{subsec:lattice_specialisation} to specialise to $n=3$.

\item[Motivating example] In this paper's usage, a section whose
   content states a claim or research direction without proving it.
   Section~\ref{sec:motivating_examples} (Boltzmann + fluid limit)
   consists of motivating examples.  Audit-table rows for motivating
   examples are by design \texttt{STUB} and stay \texttt{STUB} until
   a subsequent paper closes them.

\item[Paper~I, Paper~II, Paper~III] The first three papers in the
   A=1 DCL series.  Paper~I (\emph{Geometry First},
   doi:10.5281/zenodo.20078529): substrate priority; the bipartite
   octahedral lattice carries Lorentz structure.  Paper~II
   (\emph{Geometry Forces Physics}): the Standard-Model gauge group
   is derived from a single conservation axiom on the lattice.
   Paper~III (\emph{Tidal Ionization and the Quantum Roche Limit}):
   gradient-detuning ionization channel; quantum Roche limit;
   atomic ISCO.
\end{description}

\subsection*{Mathematical terms}

\begin{description}
\item[Automorphism algebra (per-site)] The Lie algebra of
   infinitesimal automorphisms of $\Lambda_n$ that fix a chosen
   site and preserve the bipartite structure plus the joint-tick
   rule (Paper~I).  At $n = 3$, $71$-dimensional (catalogued in
   \texttt{dcl-zoo}).  Generalisation to arbitrary $n$ is the
   subject of \S\ref{sec:automorphism_algebra}.

\item[Counting measure] The basic measure $\mu_\#$ on $\Lambda_n$
   assigning $|A|$ to a finite subset $A$.  See
   \S\ref{subsec:counting_measure}.

\item[Discrete codifferential] Formal adjoint $d^\ast$ of the
   discrete exterior derivative $d$, with respect to the inner
   product on $\Omega^k$.  See \S\ref{subsec:dec_codifferential}.

\item[Discrete connection] A choice of edge-valued data (a discrete
   1-form, in the abelian case) used to define parallel transport
   along lattice paths.  See \S\ref{subsec:dec_connections}.  Used
   in Paper~II's gauge-derivation programme.

\item[Discrete curvature] The 2-form on plaquettes that measures
   the failure of plaquette holonomy to be the identity.  See
   \S\ref{subsec:dec_connections}.

\item[Discrete exterior calculus] The discrete analogue of the
   exterior-calculus apparatus (forms, $d$, $d^\ast$, Stokes,
   Hodge) developed on the lattice cell complex.  See
   \S\ref{sec:discrete_diff_geometry}.

\item[Discrete exterior derivative] The operator $d \colon \Omega^k
   \to \Omega^{k+1}$ defined as a signed boundary sum.  Satisfies
   $d \circ d = 0$.  See \S\ref{subsec:dec_d}.

\item[Discrete-to-continuum limit] The limit $\epsilon \to 0$ of
   the rescaled lattice $\epsilon \Lambda_n$ towards
   $\mathbb{R}^n$, together with the limits of operators / function
   spaces / measures defined on the rescaled lattice.  See
   \S\ref{sec:continuum_limits}.

\item[Discrete $k$-form] A function on oriented $k$-cells of
   $\Lambda_n$.  See \S\ref{subsec:dec_forms}.

\item[Haar-analogue measure] On $\Lambda_n$, a measure invariant
   under the per-site automorphism algebra.  By
   uniqueness arguments analogous to Haar's theorem on locally
   compact groups, often coincides with the counting measure up to
   scaling.  See \S\ref{subsec:haar_analogue}.

\item[Holonomy (lattice / plaquette)] The product of discrete
   connection 1-forms along a closed lattice path.  Plaquette
   holonomy is the holonomy around an elementary plaquette; its
   failure to be the identity defines the discrete curvature.
   See \S\ref{subsec:dec_connections}.

\item[Hyperoctahedral group] $B_n$: the group of signed permutations
   of the $n$ coordinate axes.  Wreath product $\mathbb{Z}/2 \wr
   S_n$; order $2^n \, n!$.  Acts on $\Lambda_n$ preserving the
   bipartite structure.  See \S\ref{subsec:autom_generators}.

\item[Lattice spacing] The parameter $\epsilon$ used to rescale
   $\Lambda_n \subset \mathbb{Z}^n$ to $\epsilon \Lambda_n \subset
   \mathbb{R}^n$.  $\epsilon \to 0$ gives the discrete-to-continuum
   limit.  See \S\ref{subsec:limit_parameter}.

\item[Sobolev-analogue space] $\ell^{k,p}(\Lambda_n)$: discrete
   functions with discrete derivatives (via $d$) up to order $k$ in
   $\ell^p$.  Discrete analogue of $W^{k,p}(\mathbb{R}^n)$.  See
   \S\ref{subsec:sobolev_analogue}.
\end{description}

\subsection*{Kinetic-theory terms (motivating examples)}

\begin{description}
\item[Chapman-Enskog expansion] An asymptotic expansion of the
   kinetic distribution function $f$ in powers of the small
   parameter $\epsilon$, used (in the continuum case) to derive
   fluid equations from the Boltzmann equation.  See
   \S\ref{subsec:fluid_limit}.

\item[Collision invariant] A function $\phi(v)$ in the kernel of
   the (linearised) collision operator $Q$; physically: a quantity
   conserved in collisions (mass, momentum, energy).  Determines the
   conserved moments $\rho$, $u$, $e$ in the fluid limit.

\item[Collision operator] The operator $Q[f]$ on a kinetic
   distribution function $f$ encoding particle-collision dynamics.
   See \S\ref{subsec:boltzmann}.

\item[Equilibrium measure] A zero of the collision operator $Q$;
   parametrised by conserved moments $(\rho, u, e)$.  Maxwellian in
   the continuum case; Maxwellian-analogue on the lattice.  See
   \S\ref{subsec:boltzmann}.

\item[Fluid limit] The limit in which a kinetic equation
   (Boltzmann) gives way to a fluid equation (Navier-Stokes or
   Euler) under an appropriate scaling.  In this paper, framed as
   a motivating example in \S\ref{subsec:fluid_limit}.

\item[$H$-functional] The convex functional $H[f] = \sum f \log f$
   on kinetic distribution functions.  Boltzmann's $H$-theorem is
   the statement that $H$ is non-increasing under the kinetic
   dynamics.  See \S\ref{subsec:boltzmann}.

\item[$H$-theorem] Boltzmann's theorem (in the continuum) that
   $H[f]$ is non-increasing under the kinetic dynamics; equivalent
   to the second law of thermodynamics for the kinetic system.
   The lattice analogue is conjectured in \S\ref{subsec:boltzmann}
   and proof is deferred to a subsequent paper.

\item[Lattice-Boltzmann (kinetic theory)] The Boltzmann-shaped
   kinetic equation on $\Lambda_n$.  Note: distinct from the
   numerical-PDE method of the same name; here the discrete velocity
   set is a physical / framework structure, not a numerical
   discretisation.  See \S\ref{subsec:boltzmann}.

\item[Navier-Stokes equations] The continuum fluid equations
   $\partial_t \rho + \nabla \cdot (\rho u) = 0$, $\partial_t (\rho
   u) + \nabla \cdot (\rho u \otimes u + p \mathbb{I}) = \nabla
   \cdot \tau$, $\partial_t (\rho e) + \ldots$, with viscous stress
   tensor $\tau$.  Target of the fluid-limit ansatz in
   \S\ref{subsec:fluid_limit}.
\end{description}


\section{Code and Data}
\input{sections/code_and_data}

\section{Reproducibility Quickstart}
\input{sections/reproducibility}

%% ============================================================
\bibliographystyle{unsrt}
\bibliography{paper-bib/references}

\end{document}
